\chapter{Repository and tool guide}

\section{Quick start}
From the repository root:

\begin{lstlisting}[language=bash]
make test
make pdf
\end{lstlisting}

The first command runs standard-library checks. It compiles the C++ example with
warnings enabled, executes the TypeScript file with Node.js, validates SQL with
SQLite, and runs the shell check in a mocked mode. The second command runs
LaTeX, BibTeX, MakeIndex, and the final PDF build.

\section{Example inventory}
\noindent\begin{tabularx}{\textwidth}{@{}>{\raggedright\arraybackslash}p{0.3\textwidth}X@{}}
\toprule
Path & Purpose \\
\midrule
\path{examples/python/ledgerline.py} & Domain validation and exact money calculation \\
\path{examples/python/app.py} & Dependency-free health endpoint \\
\path{examples/typescript/order.ts} & Typed client-side calculation \\
\path{examples/java/RetryPolicy.java} & Bounded exponential backoff \\
\path{examples/cpp/latency_bucket.cpp} & Percentile selection over a sample \\
\path{examples/sql/schema.sql} & Relational constraints and an aggregate view \\
\path{examples/shell/healthcheck.sh} & Fail-fast health probe \\
\path{Dockerfile} and \path{compose.yaml} & Minimal container and local service orchestration \\
\bottomrule
\end{tabularx}

The examples are intentionally small. A small example can expose a rule without
requiring a framework, cloud account, or generated boilerplate. Production use
would require authentication, persistence lifecycle, request parsing, rate
limits, structured telemetry, deployment policy, and a threat model.

\section{Validation policy}
The repository distinguishes “verified here” from “instructions for another
environment.” The Python, SQL, shell, C++, and Node checks are run by the script
\path{scripts/verify_examples.py}. The Java source is dependency-free and includes a
compile command; this build environment has a Java runtime but not a compiler,
so the CI or a local JDK should compile it. The Docker setup is inspected as
source; Docker is optional in the current environment.

This distinction is important. A source file cannot be described as compiled
merely because it looks syntactically familiar. Record the toolchain and the
exact command used for verification.

\section{Updating the manuscript}
Edit chapter files rather than generated build files. Run `make pdf`, render the
PDF pages, inspect the title page, contents, code blocks, tables, and diagrams,
then run `make test`. If references change, rebuild BibTeX and the index. Keep
claims about versioned tools tied to a dated reference or a version range.
